\chapter{A Short History of an Idea}\label{ch:history}

\begin{etymology}
\textbf{Algorithm} comes from the name of the Persian mathematician
Muhammad ibn Musa al-Khwarizmi, who worked in Baghdad in the ninth century.
Latin translations of his book on Hindu--Arabic numerals began with
``\emph{Dixit Algorizmi}'', ``thus spoke al-Khwarizmi'', and his name became the
word for a step-by-step method.

\textbf{Oracle} comes from the Latin \emph{oraculum}, from \emph{orare}, ``to
speak'' or ``to pray'': the place where you asked a god a hard question. In
quantum computing an oracle is a black box that answers questions. You know what
it answers, not how.
\end{etymology}

In 1981 Richard Feynman asked why we try to simulate nature with machines that do
not work the way nature does. Quantum computing starts with that question. It was
built in layers over the next fifteen years, and each layer answered a question
the previous one had left open. After Feynman, someone asked whether there was a
universal quantum computer. Then came the first problems where such a computer
would win, the first problem anyone outside physics cared about, the objection
that it could never work, and the answer to the objection.

\section{Feynman's question (1981)}

\begin{history}{May 1981 \textperiodcentered{} Endicott House, MIT}
The conference on the Physics of Computation, organized by MIT and IBM, brought
together physicists and computer scientists at Endicott House, near Boston.
Richard Feynman gave a talk called ``Simulating Physics with Computers'',
published in 1982. Paul Benioff had already shown, in 1980, that a Turing
machine could be described with quantum mechanics, and Yuri Manin had raised a
similar point in Russian the same year. Feynman asked the question the other
way around.
\end{history}

\begin{history}{1981 \textperiodcentered{} Caltech}
Feynman's great work on quantum electrodynamics was thirty years behind him, and
had earned him the Nobel Prize in 1965. In 1981 the hard problem in particle
physics was quantum chromodynamics, the theory of quarks and gluons, and Feynman
published a paper on it that year, on Yang--Mills theory in two space dimensions.
The theory was known, but almost nothing could be computed from it. Since Kenneth
Wilson's lattice formulation of 1974, physicists had tried to extract numbers such
as the mass of the proton by simulating the theory on a grid of points, and by
1980 Michael Creutz was doing it with Monte Carlo methods on the largest computers
available. At the same time, Feynman had been talking computation for years with
his friend Edward Fredkin, who taught him computing in exchange for lessons in
quantum mechanics. When Fredkin invited him to give the opening talk at Endicott
House, Feynman at first refused.
\end{history}

The talk was about simulating physics in general, not quarks. Simulating a
molecule means describing its quantum state, and the number of numbers needed to
describe that state grows exponentially with the number of particles. Each new
electron multiplies the memory by a constant factor. A few dozen particles are
enough to exceed any classical memory.

Feynman went further than memory. A classical computer does not have to store the
whole state if it can sample it instead, and that is what Creutz's Monte Carlo
calculations did.

\begin{history}{1946--1949 \textperiodcentered{} Los Alamos}
Stanisław Ulam, recovering from an illness in 1946, wondered about the chances of
winning a game of solitaire. Counting every possible game was hopeless, so he
thought of playing a few hundred games and counting the wins. He saw at once that
the same trick could follow neutrons through the material of a bomb, a calculation
nobody could do exactly. John von Neumann turned the idea into programs for ENIAC.
Nicholas Metropolis named it \emph{Monte Carlo}, after the casino in Monaco where
Ulam's uncle used to gamble, and Metropolis and Ulam published it in 1949.
\end{history}

The idea fits in one picture. To find the area of a quarter circle, you don't need
geometry: throw random points into the square around it and count how many land
inside. The fraction that lands inside approaches the fraction of the area, which
is $\pi/4$.

\begin{figure}[htbp]
\centering
% Generated by a seeded run (numpy default_rng(1946)); do not edit by hand.
\def\mcInside{0.895/0.375, 0.671/0.491, 0.666/0.562, 0.438/0.218, 0.950/0.075, 0.536/0.721, 0.787/0.025, 0.282/0.476, 0.296/0.056, 0.695/0.051, 0.309/0.883, 0.373/0.506, 0.451/0.702, 0.380/0.325, 0.932/0.196, 0.454/0.544, 0.631/0.179, 0.255/0.108, 0.703/0.515, 0.509/0.585, 0.136/0.394, 0.794/0.409, 0.363/0.316, 0.335/0.552, 0.795/0.238, 0.322/0.721, 0.178/0.902, 0.541/0.175, 0.670/0.454, 0.105/0.555, 0.259/0.864, 0.239/0.480, 0.884/0.145, 0.761/0.547, 0.201/0.626, 0.307/0.208, 0.392/0.896, 0.562/0.532, 0.505/0.320, 0.503/0.365, 0.640/0.729, 0.132/0.086, 0.160/0.160, 0.586/0.532, 0.093/0.297, 0.262/0.265, 0.278/0.845, 0.161/0.372, 0.745/0.349, 0.386/0.661, 0.480/0.793, 0.665/0.248, 0.532/0.250, 0.452/0.248, 0.838/0.405, 0.609/0.401, 0.096/0.914, 0.452/0.312, 0.029/0.117, 0.261/0.019, 0.456/0.130, 0.093/0.935, 0.868/0.154, 0.222/0.458, 0.022/0.149, 0.443/0.588, 0.501/0.407, 0.542/0.293, 0.093/0.903, 0.399/0.058, 0.143/0.219, 0.560/0.179, 0.101/0.672, 0.910/0.083, 0.659/0.464, 0.718/0.014, 0.074/0.064, 0.872/0.323, 0.469/0.841, 0.747/0.186, 0.402/0.062, 0.731/0.423, 0.770/0.347, 0.141/0.583, 0.547/0.824, 0.093/0.673, 0.205/0.955, 0.724/0.564, 0.311/0.083, 0.293/0.661, 0.583/0.120, 0.216/0.215, 0.752/0.072, 0.588/0.228, 0.772/0.261, 0.895/0.372, 0.040/0.915, 0.354/0.489, 0.029/0.127, 0.696/0.378, 0.370/0.040, 0.897/0.025, 0.311/0.890, 0.382/0.624, 0.776/0.315, 0.718/0.313, 0.240/0.033, 0.543/0.502, 0.073/0.409, 0.343/0.032, 0.957/0.044, 0.287/0.300, 0.373/0.370, 0.164/0.809, 0.502/0.308, 0.851/0.460, 0.100/0.363, 0.501/0.161, 0.984/0.142, 0.680/0.711, 0.060/0.746, 0.327/0.598, 0.391/0.801, 0.470/0.021, 0.618/0.259, 0.543/0.754, 0.820/0.312, 0.749/0.587, 0.802/0.539, 0.517/0.479, 0.256/0.568, 0.038/0.927, 0.645/0.505, 0.566/0.302, 0.487/0.497, 0.881/0.390, 0.965/0.150, 0.853/0.184, 0.494/0.058, 0.542/0.185, 0.758/0.040, 0.229/0.312, 0.004/0.705, 0.238/0.926, 0.896/0.394, 0.488/0.483, 0.080/0.104, 0.297/0.592, 0.343/0.028, 0.795/0.375, 0.264/0.087, 0.154/0.905, 0.034/0.338, 0.069/0.571, 0.621/0.162, 0.344/0.900, 0.196/0.323, 0.623/0.339, 0.880/0.061, 0.326/0.404, 0.163/0.413, 0.574/0.565, 0.315/0.152, 0.561/0.635, 0.417/0.361, 0.489/0.468, 0.185/0.910, 0.440/0.091, 0.021/0.186, 0.377/0.251, 0.585/0.639, 0.464/0.069, 0.358/0.551, 0.255/0.272, 0.592/0.634, 0.554/0.500, 0.775/0.584, 0.189/0.518, 0.032/0.841, 0.254/0.089, 0.520/0.540, 0.022/0.209, 0.109/0.882, 0.278/0.297, 0.407/0.453, 0.624/0.051, 0.860/0.509, 0.514/0.822, 0.026/0.915, 0.419/0.822, 0.063/0.750, 0.033/0.837, 0.139/0.942, 0.578/0.794, 0.058/0.454, 0.155/0.587, 0.917/0.065, 0.455/0.153, 0.447/0.620, 0.107/0.313, 0.381/0.856, 0.355/0.461, 0.502/0.575, 0.412/0.401, 0.065/0.229, 0.246/0.625, 0.684/0.201, 0.906/0.131, 0.409/0.196, 0.486/0.124, 0.168/0.254, 0.586/0.182, 0.444/0.719, 0.627/0.220, 0.086/0.605, 0.044/0.203, 0.169/0.138, 0.057/0.173, 0.884/0.273, 0.893/0.436, 0.382/0.029, 0.318/0.382, 0.014/0.592, 0.459/0.691, 0.559/0.590, 0.388/0.334, 0.431/0.790, 0.165/0.577, 0.352/0.751, 0.257/0.180, 0.587/0.299, 0.732/0.242, 0.190/0.374, 0.385/0.849, 0.860/0.130, 0.225/0.856, 0.061/0.778, 0.142/0.110, 0.041/0.973, 0.556/0.487, 0.450/0.856, 0.011/0.202, 0.435/0.347, 0.256/0.075, 0.660/0.323, 0.028/0.867, 0.642/0.139, 0.259/0.584, 0.122/0.238, 0.216/0.693, 0.170/0.522, 0.226/0.908, 0.511/0.488, 0.193/0.205, 0.545/0.065, 0.170/0.119, 0.034/0.778, 0.506/0.836, 0.015/0.768, 0.121/0.823, 0.409/0.475, 0.446/0.657, 0.576/0.750, 0.821/0.133, 0.161/0.134, 0.677/0.001, 0.240/0.642, 0.393/0.850, 0.757/0.182, 0.463/0.304, 0.182/0.171, 0.763/0.428, 0.238/0.314, 0.650/0.470, 0.233/0.143, 0.228/0.313, 0.873/0.076, 0.444/0.176, 0.062/0.561, 0.510/0.578, 0.415/0.702, 0.069/0.399, 0.712/0.336, 0.515/0.578, 0.419/0.242, 0.650/0.085, 0.793/0.513, 0.602/0.676, 0.170/0.884, 0.576/0.096, 0.014/0.632, 0.680/0.699, 0.276/0.329, 0.719/0.617, 0.325/0.418, 0.472/0.100, 0.673/0.190, 0.851/0.320, 0.146/0.896, 0.284/0.921, 0.640/0.231, 0.490/0.171, 0.356/0.425, 0.160/0.094, 0.496/0.650, 0.484/0.348, 0.482/0.579, 0.013/0.249, 0.128/0.826, 0.506/0.654, 0.288/0.565, 0.091/0.299, 0.699/0.149, 0.784/0.284, 0.751/0.184, 0.864/0.154, 0.396/0.727, 0.875/0.334, 0.360/0.334, 0.519/0.307, 0.581/0.751, 0.098/0.634, 0.722/0.099}
\def\mcOutside{0.807/0.764, 0.907/0.525, 0.844/0.830, 0.950/0.956, 0.999/0.274, 0.811/0.932, 0.910/0.966, 0.609/0.809, 0.596/0.911, 0.864/0.546, 0.957/0.596, 0.917/0.462, 0.703/0.768, 0.996/0.968, 0.546/0.844, 0.756/0.967, 0.733/0.938, 0.810/0.742, 0.627/0.905, 0.810/0.934, 0.864/0.553, 0.829/0.582, 0.697/0.969, 0.713/0.912, 0.570/0.962, 0.735/0.949, 0.835/0.946, 0.528/0.879, 0.791/0.973, 0.931/0.938, 0.987/0.502, 0.894/0.810, 0.793/0.918, 0.624/0.844, 0.816/0.647, 0.983/0.665, 0.975/0.378, 0.668/0.847, 0.894/0.958, 0.604/0.846, 0.338/0.997, 0.760/0.982, 0.784/0.941, 0.969/0.725, 0.545/0.978, 0.872/0.694, 0.815/0.965, 0.603/0.833, 0.497/0.871, 0.729/0.764, 0.878/0.549, 0.848/0.943, 0.786/0.817, 0.637/0.869, 0.798/0.763, 0.874/0.956, 0.705/0.932, 0.170/0.991, 0.951/0.908, 0.424/0.942, 0.690/0.863, 0.612/0.846, 0.919/0.869, 0.788/0.731, 0.638/0.982, 0.944/0.804, 0.925/0.813, 0.488/0.878, 0.830/0.791, 0.391/0.959, 0.693/0.946, 0.931/0.743, 0.688/0.788, 0.773/0.943, 0.622/0.887, 0.914/0.908, 0.700/0.789}
\def\mcCurve{(10,3.2000) (12,3.0000) (14,3.1429) (16,3.2500) (19,3.3684) (22,3.4545) (26,3.3846) (30,3.3333) (35,3.4286) (41,3.3171) (48,3.2500) (56,3.2857) (65,3.2000) (76,3.2632) (89,3.3708) (104,3.3462) (122,3.4098) (142,3.3239) (166,3.2771) (194,3.2990) (227,3.2423) (265,3.2906) (310,3.2645) (363,3.2397) (424,3.2264) (495,3.1919) (579,3.1986) (677,3.1965) (791,3.2263) (925,3.2043) (1081,3.2155) (1264,3.1930) (1477,3.1984) (1727,3.1940) (2019,3.1778) (2360,3.1780) (2759,3.1751) (3225,3.1702) (3769,3.1722) (4406,3.1720) (5151,3.1722) (6021,3.1729) (7038,3.1583) (8227,3.1530) (9617,3.1552) (11242,3.1560) (13141,3.1568) (15362,3.1519) (17957,3.1546) (20991,3.1539) (24538,3.1532) (28683,3.1569) (33529,3.1610) (39194,3.1639) (45816,3.1564) (53557,3.1522) (62605,3.1527) (73182,3.1479) (85547,3.1489) (100000,3.1476)}
\def\mcShowN{400}
\def\mcShowIn{323}
\def\mcShowEst{3.23}
\def\mcFinalEst{3.1476}

\begin{tikzpicture}
  \begin{scope}[scale=4.6]
    \fill[hfConcept!6] (0,0) -- (1,0) arc (0:90:1) -- cycle;
    \draw[black!40] (0,0) rectangle (1,1);
    \draw[hfConcept, line width=0.8pt] (1,0) arc (0:90:1);
    \foreach \x/\y in \mcInside { \fill[hfConcept] (\x,\y) circle (0.008); }
    \foreach \x/\y in \mcOutside { \fill[black!45] (\x,\y) circle (0.008); }
    \node[hf/note, below=1mm, text width=44mm] at (0.5,0) {\mcShowIn{} of \mcShowN{}
      points inside\\$\pi \approx 4 \times \mcShowIn/\mcShowN = \mcShowEst$};
  \end{scope}
  \begin{semilogxaxis}[at={(66mm,0)}, width=66mm, height=52mm,
      xmin=10, xmax=100000, ymin=2.9, ymax=3.5,
      xlabel={number of points}, ylabel={estimate of $\pi$},
      ylabel style={yshift=2mm},
      tick label style={font=\sffamily\scriptsize},
      label style={font=\sffamily\small},
      axis lines=left, axis line style={black!50}]
    \addplot[hfHistory, dashed, domain=10:100000, samples=2] {pi};
    \addplot[hfConcept, line width=0.8pt, mark=none] coordinates {\mcCurve};
  \end{semilogxaxis}
\end{tikzpicture}
\caption{Monte Carlo. Left: random points in a square, counted inside and outside a
quarter circle. Right: the estimate as the number of points grows, against the true
value of $\pi$ (dashed). After 100\,000 points it reads \mcFinalEst. The error
shrinks like $1/\sqrt{N}$: each extra correct digit costs a hundred times more
points.}
\end{figure}

The same trick works on Wilson's lattice. In 1974 Kenneth Wilson replaced
continuous space and time with a grid of points a small distance $a$ apart. Quarks
live on the points, the \emph{sites}. The gluon field lives on the \emph{links}
between neighboring sites: each link carries a $3\times3$ unitary matrix $U$ that
says how a quark's color is rotated when it hops from one site to the next. Wilson
put the gluons on the links, and not on the sites, because that keeps the symmetry
of the theory exact even on a coarse grid. The smallest closed loop, a square of
four links called a \emph{plaquette}, multiplies four of those matrices. How far
the product is from the identity measures the strength of the field inside the
square, and adding it up over all plaquettes gives the \emph{action} $S$ of the
configuration.

\begin{figure}[htbp]
\centering
\begin{tikzpicture}[x=13mm, y=13mm]
  % plaquette
  \fill[hfHistory!12] (2,1) rectangle (3,2);
  \draw[hfHistory, line width=0.7pt, -{Stealth[length=1.8mm]}]
    (2.5,1.5) ++(-40:0.27) arc (-40:250:0.27);
  % links
  \foreach \x in {0,...,5} \foreach \y in {0,...,3} {
    \ifnum\x<5 \draw[hfConcept!70, line width=0.8pt] (\x,\y) -- (\x+1,\y); \fi
    \ifnum\y<3 \draw[hfConcept!70, line width=0.8pt] (\x,\y) -- (\x,\y+1); \fi
  }
  % one highlighted link
  \draw[hfConcept, line width=2pt, -{Stealth[length=2.6mm]}] (4,2) -- (4.92,2);
  \node[font=\sffamily\small, text=hfConcept, above] at (4.5,2) {$U$};
  % sites
  \foreach \x in {0,...,5} \foreach \y in {0,...,3}
    \fill[hfSource] (\x,\y) circle (2.4pt);
  % spacing
  \draw[black!60, decorate, decoration={brace, amplitude=4pt, mirror}]
    (0,-0.12) -- (1,-0.12) node[midway, below=4pt, font=\sffamily\small, text=black] {$a$};
  % axes
  \draw[black!50, -{Stealth[length=2mm]}] (-0.45,-0.7) -- (1.0,-0.7)
    node[right, hf/note] {space};
  \draw[black!50, -{Stealth[length=2mm]}] (-0.45,-0.7) -- (-0.45,0.7)
    node[above, hf/note] {time};
  % labels
  \node[hf/note, anchor=west, text width=40mm, align=left] (ls) at (5.45,3) {\textbf{site}: a point of spacetime, where the quarks live};
  \draw[black!40] (ls.west) -- (5.08,3);
  \node[hf/note, anchor=west, text width=40mm, align=left, text=hfConcept] (ll) at (5.45,2) {\textbf{link}: the gluon field, a $3\times3$ matrix $U$};
  \draw[hfConcept!50] (ll.west) -- (5.0,2);
  \node[hf/note, anchor=west, text width=40mm, align=left, text=hfHistory] (lp) at (5.45,1) {\textbf{plaquette}: four links around a square; their product measures the field};
  \draw[hfHistory!60] (lp.west) -- (3.05,1.5);
\end{tikzpicture}
\caption{A slice of Wilson's lattice, two of its four dimensions. One
\emph{configuration} is a choice of matrix for every link. Monte Carlo samples
configurations instead of summing over all of them.}
\end{figure}

A configuration is one choice of matrix for every link, and the theory says to sum
over all of them. A modest lattice of $32^4$ sites has 4\,194\,304 links, and each
$3\times3$ matrix of this kind has 8 free real parameters: the sum is an integral
over more than 33 million variables. No computer can do that integral directly.
Instead, you sample configurations at random, each with probability proportional to
$e^{-S}$, and average over the samples, exactly like the points in the square.

It works because $e^{-S}$ is a positive number, so it behaves like a probability.
And it is positive only thanks to a trick: Wilson's lattice runs in \emph{imaginary}
time. In real time, quantum mechanics weighs each configuration with $e^{iS}$, a
complex number that turns around the circle as $S$ changes. Replacing $t$ with $-i\tau$
turns that rotating weight into a positive one. For a proton at rest, which does not
change in time, nothing is lost. For anything that moves, like a collision, the
trick is not available.

Feynman showed that this cannot work for quantum mechanics in general. Quantum
amplitudes are not positive numbers; they carry signs and phases, and they cancel
(Chapter~\ref{ch:qubit}). When the weights can be negative, random samples mostly
cancel each other and the noise swamps the answer. Physicists later named this the
\emph{sign problem}. The correlations that quantum mechanics predicts, the ones John
Bell had turned into an inequality in 1964, cannot be reproduced by any local
classical process running on ordinary probabilities. ``Nature isn't classical, dammit,'' he
said, ``and if you want to make a simulation of nature, you'd better make it quantum
mechanical.''

Feynman's proposal was simple: use a quantum system you can control to imitate
the quantum system you want to study. Prepare it, let it evolve according to the
dynamics of interest, measure at the end. The physics does the work instead of
the processor. At that point the idea was simulation, not general-purpose
computing.

\section{A universal machine (1985)}

In 1985, David Deutsch, at Oxford, asked the general question: is there a model
of quantum computation that is \emph{universal}, able to simulate any physical
system efficiently? He defined the \emph{quantum Turing machine} and argued that
it was. In the same paper he gave the first problem where a quantum machine does
better than a classical one: deciding whether a function of one bit is constant
or balanced. A classical machine has to ask the function twice. Deutsch's
version asked once and got the answer half of the time; the version that always
works came in 1998.

The problem is artificial and the gain is tiny, but the separation is clean. This
is where quantum computing becomes a model of computation, not just a way of
simulating physics.

\section{The algorithms (1992--1996)}

\begin{history}{1992--1993 \textperiodcentered{} Oxford and Berkeley}
In 1992, Deutsch and Richard Jozsa extended the 1985 problem to functions of $n$
bits: the quantum algorithm needs two queries, while a deterministic classical
algorithm needs, in the worst case, $2^{n-1}+1$. Richard Cleve, Artur Ekert,
Chiara Macchiavello and Michele Mosca brought it down to one query in 1998,
which is the form taught today. In 1993,
Ethan Bernstein and Umesh Vazirani put the quantum Turing machine on rigorous
ground and defined BQP, the class of problems a quantum computer solves
efficiently (Chapter~\ref{ch:turing}).
\end{history}

Deutsch--Jozsa is a toy. A classical algorithm that is allowed to guess at random
solves it with high probability in a handful of queries, and nobody needs to know
whether a function is constant. Its value was to show that the separation could
be large, not to solve anything.

\begin{history}{1994 \textperiodcentered{} STOC, FOCS and Bell Labs}
Daniel Simon found a problem with an exponential separation, even against
randomized classical algorithms: finding a hidden period in a function defined
on the corners of an $n$-dimensional cube. His method was, in essence, a Fourier
transform over those corners. He sent the paper to STOC, one of the two main
conferences in theoretical computer science, and it was rejected. Peter Shor was
on the program committee. He later said he argued for accepting it, though not
hard enough.

Shor had seen what Simon's paper meant. Fourier transforms find periods; the
discrete logarithm problem behind Diffie--Hellman is a period problem on a torus
instead of a cube. He solved the discrete logarithm first and presented it in an
internal seminar at Bell Labs in the spring of 1994. A few days later he had
factoring.
By the time Umesh Vazirani called him, the rumor had already turned ``discrete
log'' into ``factoring''. Simon's paper and Shor's were both accepted at FOCS, in
November 1994.
\end{history}

Shor's insight was that factoring and the discrete logarithm are period-finding
problems, the same kind as Simon's, only on a cyclic group instead of a cube.
Chapter~\ref{ch:factoring} shows the reduction. The quantum Fourier transform that finds the
period is the subject of Chapter~\ref{ch:fourier}.

\begin{history}{1994--1996 \textperiodcentered{} Santa Fe and Bell Labs}
In May 1994, at a workshop at the Santa Fe Institute, Rolf Landauer of IBM
raised the objection that seemed to end the story: a quantum computer cannot
correct its errors. Every classical technique for fixing errors involves copying
something, and quantum states cannot be copied. Without correction, a
computation of a billion steps needs each step to be accurate to one part in a
billion. In 1995 Shor answered with a code that protects one qubit using nine.
In 1996, Robert Calderbank and Shor, and independently Andrew Steane, generalized
it into a whole family, the CSS codes. Part~\ref{part:qec} tells this story in detail.
\end{history}

In 1996, Lov Grover, also at Bell Labs, published a quantum algorithm for
unstructured search: finding the one marked item among $N$ takes about
$\sqrt{N}$ queries instead of $N$. The gain is only quadratic, much less dramatic
than Shor's. But it applies to almost anything that can be phrased as a search,
which made it the second pillar of the field.

\section{What the sequence shows}

By 1996 every piece of the argument was on the table. Feynman showed that
classical machines struggle with quantum physics. Deutsch showed there is a
universal model that does not. Simon showed that the Fourier transform is the
mechanism behind exponential gains. Shor showed that one of those gains breaks
the cryptography of the internet, which turned an academic curiosity into
something governments fund. And error correction showed that the model can, in
principle, survive noise.

What did not exist in 1996 was a machine. The thirty years since then are the
story of Parts~\ref{part:hardware} and~\ref{part:qec}.
